% 2.6 — La surveillance à l'exécution
\section{La surveillance à l'exécution}
\label{sec:execution}

Ce qu'aucune lecture ne peut prédire --- un script Python, un \cmd{make}, une commande
construite à l'exécution --- ne peut être vu qu'en \emph{regardant} la commande travailler.
Les systèmes d'exploitation offrent plusieurs points d'observation.

\begin{description}[style=nextline,leftmargin=0.6cm,font=\sffamily\bfseries]
  \item[ptrace (Linux)]
    L'appel système \cmd{ptrace} permet à un processus traceur « d'observer et de
    contrôler l'exécution d'un autre processus, et d'examiner et modifier sa mémoire et ses
    registres » ; il sert d'abord aux débogueurs et au traçage des appels
    système~\cite{man-ptrace}. Un traceur peut s'arrêter à l'entrée de chaque appel système
    du processus suivi --- donc \emph{avant} qu'un \cmd{unlink} ne s'exécute.
  \item[seccomp avec notification (Linux)]
    Un filtre seccomp peut déléguer la décision sur un appel système à un processus
    superviseur en espace utilisateur, plus rapidement qu'avec ptrace puisque seuls les
    appels filtrés sont notifiés. La page de manuel avertit toutefois que ce mécanisme
    « n'est pas destiné à implémenter une politique de sécurité », notamment à cause d'une
    situation de concurrence entre la vérification et l'usage des
    arguments~\cite{man-seccomp-unotify}.
  \item[Landlock (Linux)]
    Landlock permet à un processus, même non privilégié, de restreindre ses propres droits
    d'accès~\cite{landlock}. C'est un outil de \emph{restriction} --- celui des bacs à sable ---
    pas d'observation ni de copie.
  \item[FUSE (Linux)]
    Un système de fichiers en espace utilisateur, dont « les données et métadonnées sont
    fournies par un processus ordinaire »~\cite{fuse-kernel}. Placé sous le projet, il voit
    chaque opération sur ses fichiers.
  \item[WinFsp (Windows)]
    L'équivalent de FUSE pour Windows : il permet d'écrire un système de fichiers comme un
    programme en mode utilisateur, sans programmation noyau~\cite{winfsp}. Windows n'ayant
    pas de ptrace, c'est l'endroit où toutes les routes (Win32, ntdll, appels système
    directs) convergent avant d'atteindre les fichiers.
  \item[Pilotes de filtre noyau (Windows)]
    Les \textit{minifilters} interceptent les opérations du système de fichiers dans le
    noyau. Ils sont puissants mais exigent un pilote signé et une programmation noyau.
\end{description}

\begin{pourquoi}[Observer n'est pas comprendre, mais c'est un plancher]
Un filet d'exécution ne sait pas \emph{pourquoi} un fichier est détruit ; il sait
seulement \emph{qu'il va l'être}, au moment où l'appel système arrive. C'est exactement ce
qui manque pour le code que \demo{} ne lit pas : on ne peut pas obscurcir un appel
système --- \cmd{rm}, quelle que soit la façon dont il est écrit, devient un \cmd{unlink}
sur un chemin résolu par le noyau. Le filet ne voit cependant que les fichiers locaux : une
base distante ou un dépôt distant lui sont invisibles.
\end{pourquoi}
